\subsection{Secondary analysis: intensity-only agent}\label{sec:secondary}
The component study (Section~\ref{sec:component}) suggested that restricting AA-DQN to the intensity actions \{\textsc{easy}, \textsc{hard}, \textsc{break}\} improves performance. We therefore trained seven additional seeds of this variant and evaluated it with the same ten-seed protocol as the primary agent. Because the variant was selected after a look at test-set results from three seeds, we report it as a secondary, hypothesis-generating analysis and keep the pre-specified hybrid agent as the primary result.

With ten seeds, the intensity-only agent achieves a return of \AADQNThreeactReturn{}~$\pm$~\AADQNThreeactReturnSD{} and an on-time rate of \AADQNThreeactOnTimePct\%. Relative to the best rule, \bestHeur, the paired differences are $\SecRetSPTThreshold$ in return ([\SecRetSPTThresholdLo, \SecRetSPTThresholdHi], $\pSecRetSPTThreshold$), $\SecOnSPTThreshold$ pp in on-time rate ([\SecOnSPTThresholdLo, \SecOnSPTThresholdHi], $\pSecOnSPTThreshold$) and $\SecWOnSPTThreshold$ pp in difficulty-weighted on-time rate ([\SecWOnSPTThresholdLo, \SecWOnSPTThresholdHi]). Relative to the primary hybrid agent, the differences are $\SecRetAADQN$ in return ([\SecRetAADQNLo, \SecRetAADQNHi]) and $\SecOnAADQN$ pp in on-time rate. The mean fatigue is \AADQNThreeactFatigue{} and the break share \AADQNThreeactBreakPct\%. Taken together, the primary and secondary analyses show that a memory-augmented DQN reaches, and in its compact variant exceeds, the performance of the best of nine tuned rules. They also show that a compact, cognitively meaningful action space is preferable to a larger rule-selection space for this problem.
